\section{One actual original line family for any finite family of numerical classes}
\label{sec:native-numerical-finite-class-families}

\begin{theorem}[Simultaneous original numerical-class realization {\cite[\S1.2]{BJ22}}]
\label{thm:native-numerical-finite-class-families}
\lean{EnvelopeProperty.NativeRealNumericalClasses.exists_common_finite_line_family, EnvelopeProperty.NativeRealNumericalClasses.exists_common_generated_line_family}
\leanok
\uses{thm:native-numerical-line-family-maps, thm:native-numerical-generated-line-sums}
Every finite family of classes in the same actual original real
numerical quotient has a single internally constructed finite
family of actual original invertible lines, indexed by some
$\operatorname{Fin}(n)$, whose whole coefficient map represents
all the given classes. One actual original common invertible
denominator $T$, genuine finite global generators of $T$, and
genuine finite global generators of all the original numerators
can be produced simultaneously. Thus every given class has its
exact real expression by this same finite family of true generated
numerator-minus-common-denominator classes.
\end{theorem}

\begin{proof}
For each of the finitely many given classes, take the actual finite
set of original invertible lines in its internally produced finite
expression. Take the genuine finite union of these finite sets.
Extend each original coefficient function to this same finite union
by zero outside its own support. The finite-sum subset identity
shows that each original class retains its exact value on this
common family. Reindex the actual finite union by its genuine
equivalence with $\operatorname{Fin}(n)$, preserving every original
line and all coefficients for every given class simultaneously.

Apply the proved original finite-line-family global-generation
producer once to this same common family. It constructs one
actual invertible denominator $T$, true finite generators of it
and of every numerator $L_i\otimes T$, and gives the exact
numerical difference formula on the entire original coefficient
space. Evaluate that single identity on each of the already
constructed coefficient vectors. The denominator, numerators and
generators are therefore common to all given classes; none is
supplied as a desired realization or generation certificate. This
finite-family result makes no finite-dimensionality assertion
about the whole numerical space and uses no ample-cone or model
psh numerical-invariance comparison.
\end{proof}
